% Part: normal-modal-logic
% Chapter: sequent-calculus
% Section: proofs-in-K

\documentclass[../../../include/open-logic-section]{subfiles}

\begin{document}

\olfileid{nml}{seq}{prk}

\olsection{\Log{K} માટે સિક્વન્ટ \usetoken{P}{derivation}}

\iftag{prvBox}{

\begin{ex}
  $\Proves (\Box!A \land \Box!B) \lif \Box (!A \land !B)$
  બતાવતું સિક્વન્ટ કલનનું !!{derivation} આપીએ છીએ.
  \begin{prooftree}
    \Axiom$!A \fCenter !A$
    \doubleLine
    \UnaryInf$!B, !A \fCenter !A$   
    \Axiom$!B \fCenter !B$
    \doubleLine
    \UnaryInf$!B, !A \fCenter !B$   
    \RightLabel{\RightR{\land}}
    \BinaryInf$!B, !A \fCenter !A \land !B$    
    \RightLabel{$\Box$}
    \UnaryInf$\Box!B, \Box!A \fCenter \Box
    (!A \land !B)$
    \RightLabel{\LeftR{\land}}
    \UnaryInf$\Box!A \land \Box!B, \Box!A \fCenter \Box
    (!A \land !B)$
    \RightLabel{\LeftR{\Exchange}}
    \UnaryInf$\Box!A, \Box!A \land \Box!B \fCenter \Box
    (!A \land !B)$
    \RightLabel{\LeftR{\land}}
    \UnaryInf$\Box!A \land \Box!B, \Box!A \land \Box!B \fCenter \Box (!A \land !B)$
    \RightLabel{\LeftR{\Contraction}}
    \UnaryInf$\Box!A \land \Box!B \fCenter \Box (!A \land !B)$
    \RightLabel{\RightR{\lif}}
    \UnaryInf$\fCenter (\Box!A \land \Box!B)
    \lif \Box (!A \land !B)$
  \end{prooftree}
\end{ex}
}{}

\iftag{prvDiamond}{
  \begin{ex}
    $\Proves \Diamond(!A \lor !B) \lif
    (\Diamond !A \lor \Diamond!B)$ બતાવતું સિક્વન્ટ કલનનું
    !!{derivation} આપીએ છીએ.
    \begin{prooftree}
      \Axiom$!A \fCenter !A$
      \doubleLine
      \UnaryInf$!A \fCenter !A, !B$   
      \Axiom$!B \fCenter !B$
      \doubleLine
      \UnaryInf$!B \fCenter !A, !B$   
      \RightLabel{\LeftR{\lor}}
      \BinaryInf$!A \lor !B \fCenter !A, !B$    
      \RightLabel{$\Diamond$}
      \UnaryInf$\Diamond(!A \lor !B) \fCenter 
      \Diamond !A, \Diamond!B$
      \RightLabel{\RightR{\lor}}
      \UnaryInf$\Diamond(!A \lor !B) \fCenter 
      \Diamond !A, \Diamond !A \lor \Diamond!B$
      \RightLabel{\RightR{\Exchange}}
      \UnaryInf$\Diamond(!A \lor !B) \fCenter \Diamond !A \lor \Diamond!B, \Diamond !A$
      \RightLabel{\RightR{\lor}}
      \UnaryInf$\Diamond(!A \lor !B) \fCenter 
      \Diamond !A \lor \Diamond!B, \Diamond !A \lor \Diamond!B$
      \RightLabel{\RightR{\Contraction}}
      \UnaryInf$\Diamond(!A
      \lor !B) \fCenter \Diamond !A \lor \Diamond!B$
      \RightLabel{\RightR{\lif}}
      \UnaryInf$\fCenter \Diamond(!A
      \lor !B) \lif (\Diamond !A \lor \Diamond!B)$
    \end{prooftree}
  \end{ex}
}{}

\iftag{notprvBox,notprvDiamond}{}{
  અહીં~\Dual નું !!a{derivation} છે.
  \sourcecorrection{OLMD-075}{નીચેના વૃક્ષની પ્રથમ
  $!A \fCenter !A$ પંક્તિમાંથી $\lnot !A,!A \fCenter$
  મેળવવા ડાબી નિષેધનો નિયમ જોઈએ; મૂળમાં તેને જમણી
  નિષેધનો નિયમ કહ્યો છે. નિયમનું લેબલ સુધાર્યું છે.}
  \begin{prooftree}
    \Axiom$!A \fCenter !A$
    \RightLabel{\LeftR{\lnot}}
    \UnaryInf$\lnot !A, !A \fCenter $
    \RightLabel{$\Diamond$}
    \UnaryInf$\Diamond\lnot !A, \Box !A \fCenter $
    \RightLabel{\RightR{\lnot}}
    \UnaryInf$\Box !A \fCenter \lnot\Diamond\lnot !A$
    \RightLabel{\RightR{\lif}}
    \UnaryInf$\fCenter \Box !A \lif \lnot\Diamond\lnot !A$
    \Axiom$!A \fCenter !A$
    \RightLabel{\RightR{\lnot}}
    \UnaryInf$\fCenter !A, \lnot !A $
    \RightLabel{\RightR{\Exchange}}
    \UnaryInf$\fCenter \lnot !A, !A $
    \RightLabel{$\Box$}
    \UnaryInf$\fCenter \Diamond\lnot !A, \Box !A$
    \RightLabel{\RightR{\Exchange}}
    \UnaryInf$\fCenter \Box !A, \Diamond\lnot !A$
    \RightLabel{\RightR{\lnot}}
    \UnaryInf$\lnot\Diamond\lnot !A \fCenter \Box !A$
    \RightLabel{\RightR{\lif}}
    \UnaryInf$\fCenter \lnot \Diamond \lnot !A \lif \Box !A$
    \RightLabel{\RightR{\land}}
    \BinaryInf$\fCenter \Box !A \liff \lnot\Diamond\lnot !A$ 
  \end{prooftree} 

% અભ્યાસપ્રશ્નો મુલતવી હોય ત્યારે આ ચાલતું નથી.
% \begin{prob}
%     $\Log{K}$માં $\Diamond!A \liff \lnot\Box\lnot !A$નું
%     !!a{derivation} આપો.
% \end{prob}
}

\begin{prob}
  $\Log{K}$માં નીચેના !!{formula}s માટે સિક્વન્ટ કલનનાં
  પ્રમાણો શોધો:
  \begin{enumerate}
    \item $\Box \lnot p \lif \Box(p \lif q)$
    \item $(\Box p \lor \Box q) \lif \Box(p \lor q)$
    \item $\Diamond p \lif \Diamond(p \lor q)$
    \item $\Box(p \land q) \lif \Box p$
  \end{enumerate}
\end{prob}

\end{document}
